\documentclass[12pt,titlepage,a4paper]{article}

\usepackage{xcolor}
\usepackage{amsmath}
\usepackage{titlesec}
\usepackage[portuguese]{babel}
\usepackage[colorlinks=true,allcolors=black]{hyperref}

\title{Data Path e Control Path}
\author{Fernando de Lima Gonçalves}
\date{Belo Horizonte\\2026}

\newcommand{\sectionbreak}{\clearpage}

\begin{document}
\maketitle
\tableofcontents

\section{Introdução}
\subsection{Data Path}
Este é o "músculo" do processador.
Elementos que executam operações (como a \texttt{ULA}, somadores do PC e para calcular um salto relativo).

\subsection{Control Path}
Este é o "cérebro" da CPU.
A partir de alguns bits do \texttt{machine code} da instrução, algumas partes do processador são acionadas, enquanto outras, não.
\texttt{MUX} e \texttt{DEMUX} são \textbf{também} elementos do control path.
Alguns bits da instrução indicam a seleção do mux (ex. jar, ativa o salto relativo a uma posição na memória).

\section{Desenvolvimento}
Assim como não existe mente sem corpo e corpo sem mente, \texttt{control} e \texttt{data path} se complementam.

Um gargalo temporal do processador é descrito por operações de leitura e escrita na memória de dados,
\textbf{muito mais lentas}. Isso faz com que, a instrução em memória demore muito mais que instruções de acesso aos registradores.
Instruções têm tempos de execução variáveis, no modelo Monociclo (cada execução é executada em 1 ciclo de clock) e, portanto,
o clock é definido, em sua frequência, pelo tempo de execução da instrução mais lenta representado por:
\begin{center}
    \large
    \[
        \boxed{f_{clock} = \frac{1}{T}}
    \]
\end{center}

\end{document}
